\linespread{1.18}\selectfont
\setlength{\parskip}{.25em}
\setcounter{figure}{0}
\renewcommand{\thefigure}{引.\arabic{figure}}
\chapter*{引入：记录还停在从前}
\phantomsection\label{front:introduction}
\addcontentsline{toc}{chapter}{引入：记录还停在从前}
\begin{center}\BookEpigraph{introduction}\end{center}


晚班交接时，实验员 Alice 接到一份不能重新取得的样本。上一班的 Bob 已经离开，操作单夹在样本旁边，前面的检查项目都有签名，只有“加入试剂”一栏空着。按原定安排，下一步应该由她继续。负责人在工作群里问进度，希望这项任务今晚能完成。

Alice 给 Bob 发了条消息：“试剂加过了吗？”等了一会儿，没有回复。她核对了样本编号，没拿错。可 Bob 会不会已经加过，只是忘了登记？她把试剂放在一边，在交接单上记下情况，请负责人先暂停这一步，等问清楚再继续。

负责人起初只看到“加入试剂”一栏空白，提醒她按流程完成；等 Alice 说明样本无法重取，才同意保留现场。此时 Bob 还没回复，这一栏也无法说明试剂是否已经加入。

\begin{figure}[htbp]
\centering\includegraphics[width=\textwidth]{assets/illustrations/lab_handover.png}
\caption{交接时，样本和操作单一起留了下来；它们未必记录着同一个时刻。}
\end{figure}

Bob 离开前确实加过试剂。他放好器具，正准备登记，隔壁工位有人急着请他协助处理一件事务。他先去帮忙，回来时又接到临时通知，随后离开了实验室。登记这一笔一直没补上。样本发生了变化，操作单却保持着动作之前的样子。

Bob 如果还没动手，也会留下同样一张“加入试剂”栏空白的操作单。Alice 仅凭这一栏，分不清试剂到底加过没有（图\ref{fig:lab-fork}）。

\begin{figure}[htbp]\centering
\begin{tikzpicture}[x=1cm,y=1cm]
\node[anchor=east] at (0,1) {过去甲};
\node[anchor=east] at (0,-.5) {过去乙};
\draw[flow] (.2,1)--(4,1);\draw[flow] (.2,-.5)--(4,-.5);
\node[event] at (1,1) {};\node[above] at (1,1.1) {已加入试剂};
\node[above] at (1,-.4) {尚未加入试剂};
\node[align=center,anchor=west] at (4.5,.25) {交接时都看到\\同一步骤栏空白};
\draw[flow,draw=Accent] (4,1)--(4.5,.4);\draw[flow,draw=Accent] (4,-.5)--(4.5,.1);
\end{tikzpicture}
\caption{相同的空白记录，可以来自两段不同的过去。}\label{fig:lab-fork}
\end{figure}


把规定改成“先打勾，再动手”又会怎样？Bob 可能刚打完勾就被叫走，试剂还没加。Alice 接班后若照单跳过这一步，就会漏加。沿用“先做再记”的顺序，Bob 加完试剂后没登记，Alice 接班时再加一次，就会重复。图\ref{fig:gaps}列出了这两种中断，以及照单继续可能带来的后果。

\begin{figure}[htbp]\centering
\begin{tikzpicture}[x=1cm,y=1cm,
  stage/.style={draw=Ink!25,fill=Paper!55,rounded corners=2pt,
    minimum height=2.65cm,inner sep=0pt},
  consequence/.style={stage,draw=Accent!50,fill=Accent!5}]
\node[font=\small\bfseries] at (2.1,6.6) {Bob 在这里中断};
\node[font=\small\bfseries] at (7.45,6.6) {交接时留下的状态};
\node[font=\small\bfseries] at (12.55,6.6) {若 Alice 直接照单继续};

\foreach \y in {4.85,1.4}{
  \node[stage,minimum width=4.15cm] at (2.1,\y) {};
  \node[stage,minimum width=3.8cm] at (7.45,\y) {};
  \node[consequence,minimum width=4.5cm] at (12.55,\y) {};
  \draw[flow] (4.25,\y)--(5.45,\y);
  \draw[flow] (9.45,\y)--(10.2,\y);
}

% 上排：只有登记发生了，操作在中断时尚未执行。
\node[font=\small\bfseries] at (2.1,5.78) {先记再做};
\node at (2.1,5.28) {勾选“已完成”};
\draw[draw=Ink,line width=.7pt] (2.1,5.02)--(2.1,4.85);
\draw[Accent,line width=1.2pt] (1.98,4.85)--(2.22,4.61)
  (1.98,4.61)--(2.22,4.85);
\node[anchor=west,text=Accent,font=\footnotesize] at (2.45,4.73) {此处中断};
\draw[draw=Muted,dashed] (2.1,4.55)--(2.1,4.34);
\node[text=Muted] at (2.1,4.03) {加试剂：尚未执行};
\node at (7.45,5.2) {操作单：已勾选};
\node at (7.45,4.5) {样本：未加试剂};
\node at (12.55,5.4) {Alice 跳过这一步};
\draw[flow,draw=Accent] (12.55,5.12)--(12.55,4.79);
\node[text=Accent,font=\large\bfseries] at (12.55,4.35) {漏加试剂};

% 下排：只有加试剂发生了，登记在中断时尚未完成。
\node[font=\small\bfseries] at (2.1,2.33) {先做再记};
\node at (2.1,1.83) {加入试剂};
\draw[draw=Ink,line width=.7pt] (2.1,1.57)--(2.1,1.4);
\draw[Accent,line width=1.2pt] (1.98,1.4)--(2.22,1.16)
  (1.98,1.16)--(2.22,1.4);
\node[anchor=west,text=Accent,font=\footnotesize] at (2.45,1.28) {此处中断};
\draw[draw=Muted,dashed] (2.1,1.1)--(2.1,.89);
\node[text=Muted] at (2.1,.58) {登记：尚未完成};
\node at (7.45,1.75) {该步骤栏：仍空白};
\node at (7.45,1.05) {样本：已加过试剂};
\node at (12.55,1.95) {Alice 再加一次};
\draw[flow,draw=Accent] (12.55,1.67)--(12.55,1.34);
\node[text=Accent,font=\large\bfseries] at (12.55,.9) {重复加入试剂};
\end{tikzpicture}
\caption{中断后，操作单与样本留下不同状态。若接班人直接照单继续，可能漏加，也可能重复加入试剂。}\label{fig:gaps}
\end{figure}


交接时可以核对操作单和签名，漏记的步骤还得另外查证。检查样本能否查出试剂已经加入，要看实验本身有没有可靠、无损的检测办法。Bob 一直联系不上时，也得有人决定继续等待，还是终止这次实验。

后来 Bob 回了消息，确认试剂已经加过。负责人把说明补进交接记录，再安排后续实验。要是 Bob 一直没回，又没有其他可靠证据，这一步就仍待核查。迟来的消息也可能添乱：如果他说“我现在去加”，Alice 还得先问清楚他指的是哪一步，并告诉他这份样本已经交接。

即使只指定一个人负责，也不能让时间停住。那个人会下班，可能离开岗位，也可能在接手者已经获得授权后才收到旧通知。交接需要说明从什么时候起由谁决定，旧的安排又凭什么失效。否则，两个人都可能按照各自收到的要求尽责，最后却对同一件事各做了一次。

同样，靠“样本编号加试剂名称”来辨认动作也不够。这份样本在后面的另一道工序里，可能还要合法地加入同一种试剂。眼前两次相似的动作，究竟是对同一步的补做，还是原本就应当分别完成的两步，要由任务本身来区分。只看材料相同，会把合法的新工作挡住；只看谁来执行，又可能把交接后的重做认成新工作。

我们先把样本留在这里。它提醒人注意一条很窄的缝隙：世界已经改变，能够传给下一个人的记录却还没有改变。计算机能把一些写入放在同一个事务中，但系统也会遇到无法一并提交的外部动作。更出人意料的是，即使记录当时齐全，恢复旧备份也可能让系统重新落回那条缝隙。

\section*{一份较早的世界}

玩家 Bob 把角色从甲地域迁往乙地。迁移完成后，他在乙地登录，使用了一批材料升级装备，又把一件旧装备转交给同伴 Alice。Alice 带着装备参加了之后的活动。对两个人来说，迁移已经过去，新的游戏经历建立在它成功的基础上。

过了一段时间，甲地域因故障恢复到迁移前的一份备份。甲地的备份包含旧角色、当时尚未消耗的材料、仍在角色名下的装备，以及一项待执行的迁移任务；乙地和交易相关的历史没有随它一起退回。

恢复程序读取甲地，看到的每一项都能对得上。角色在，物品在，任务也在。如果只检查这份快照内部的关系，它甚至显得很整齐。程序于是准备再次执行迁移。它并不知道，快照以外已经有一个角色在乙地活动过，旧材料的一部分已经消耗，那件装备也有了新的持有人。

\begin{figure}[htbp]
\centering\includegraphics[width=\textwidth]{assets/illustrations/world_rollback.png}
\caption{甲地恢复到较早的时刻，乙地与其他玩家的经历仍在继续。}
\end{figure}

值班工程师先暂停了重放。他在乙地找到了角色，也看到了迁移后的使用记录。甲地恢复出了迁移前的状态，乙地保留着迁移后的活动。恢复程序准备重新执行的那项任务，早已在乙地完成过。图\ref{fig:regions}中，甲地退回了旧快照，乙地和其他玩家的活动仍在继续。

\begin{figure}[htbp]\centering
\begin{tikzpicture}[x=1cm,y=1cm]
\node[anchor=east] at (0,1.5) {甲地};\node[anchor=east] at (0,0) {乙地};\node[anchor=east] at (0,-1.5) {Alice};
\foreach \y in {1.5,0,-1.5}{\draw[flow] (.3,\y)--(12,\y);}
\node[event] at (1,1.5) {};\node[above] at (1,1.6) {备份};
\draw[flow,draw=Accent] (3,1.5)--(3,0);\node[anchor=west] at (3.1,.8) {迁移};
\node[event] at (5.7,0) {};\node[above] at (5.7,.1) {消耗材料};
\draw[flow,draw=Accent] (8,0)--(8,-1.5);\node[anchor=west] at (8.1,-.8) {转交装备};
\node[event] at (10.8,1.5) {};\node[above] at (10.8,1.6) {恢复};
\draw[flow,draw=Accent,dashed] (10.8,1.6) to[bend right=28] node[above] {仅甲地回到旧快照} (1,1.6);
\node[below] at (10,-1.6) {外部历史继续};
\end{tikzpicture}
\caption{恢复边界只覆盖甲地，迁移后的使用与交易留在边界之外。}\label{fig:regions}
\end{figure}


如果把甲地的数据再覆盖到乙地，角色也许会回到迁移时的样子，后来消耗材料换来的变化就可能被抹掉。若选择合并，两边都保留下来，那些已经消耗的材料可能重新出现，旧装备也可能再得到一份。按角色名字匹配，或者给这次传输再生成一个新的请求编号，都没有回答这批物品为什么又有资格进入乙地。

直接拒绝也需要依据。乙地此刻有角色，并不自动证明每一个待迁移任务都属于已完成的那一次。正常业务还可能存在迁回、再次迁出，或者一次迁移尚未完成便发生故障。工程师需要知道，旧任务对应的究竟是哪一次迁移，哪一方仍有权推进它，以及哪些证据在甲地回退之后依然可信。工程师可以按任务编号查找对应的完成记录，确认旧任务是否已经处理过。

Bob 在客服页面补充说，旧装备已经给了 Alice。这个细节使处理范围又扩大了一步。即便能把 Bob 的角色完整地退回迁移前，也不能假定 Alice 同意撤销之后的活动。装备可能已经影响过一次战斗，或继续进入下一笔交易。恢复甲地时没有征询这些参与者，事后也不能只靠改回甲地的一条记录，把外面的变化当作从未发生。

\begin{figure}[htbp]\centering
\begin{tikzpicture}[x=1cm,y=1cm]
\node[align=center] (old) at (0,1) {甲地旧快照\\装备归 Bob};
\node[align=center] (new) at (5,1) {乙地迁移后\\Bob 持有装备};
\node[align=center] (other) at (10,1) {后续交易\\Alice 持有装备};
\draw[flow] (old)--node[above]{迁移}(new);\draw[flow] (new)--node[above]{转交}(other);
\draw[Accent,dashed,line width=.8pt] (-1.5,0)--(1.5,0);
\node[align=center,Accent] at (0,-.7) {恢复重新带回\\旧归属声明};
\draw[flow,draw=Accent,dashed] (1.5,-.6)--(9,-.6)--(9,.25);
\node[fill=white,align=center] at (5.4,-.6) {若重新授权，将影响已完成的交易};
\end{tikzpicture}
\caption{同一件装备的归属延伸到了第三方，旧声明不能只在甲地内部解释。}\label{fig:ownership}
\end{figure}


图\ref{fig:ownership}只画出一笔转交，已经足以说明边界：旧快照里写着“归 Bob”，外部较新的记录写着“已交给 Alice”。两条记录在各自的时间点都可能正确。系统需要保护同一件物品沿时间延续的归属关系，防止旧快照让 Bob 再次获得处置权。恢复还会影响谁能继续使用装备，以及已经完成的活动如何处理。

值班工程师继续暂停重放，核对迁移与交易记录。甲地启动以后，还要逐项检查旧任务是否已经执行、装备是否已经转交，再决定哪些工作可以继续。这些检查依赖故障前留下的记录，恢复程序需要从可靠的位置把它们读回来。

证据也不会永远免费存在。一次迁移过去很久之后，系统也许已经清理了任务记录；备份却可能比清理更早。旧任务回来时，接收方找不到成功记录，究竟是从未处理，还是处理过而后遗忘了？实验室里的空白栏又出现了，只不过这次空白来自时间中的另一个位置。保留多久、允许恢复多旧的备份、超过边界后怎样处置，都开始影响原先那个“重复执行不会多出效果”的承诺。

两个故事到这里都没有一个可贴在所有系统上的答案。接班人需要在证据不足时控制下一步动作，值班工程师需要阻止旧状态重新支配已经向前发展的历史。他们面对的对象不同，却都不能只凭眼前的一次调用或一张局部记录，决定过去是否发生、现在是否应当再做。

本书将把这些复杂处境缩小到一个可运行的订单与付款系统。缩小是为了能亲手安排每一步的先后，能在进程停止后检查证据，也能清楚地说出某一次修改究竟解决了哪一种失败。我们先让系统正常完成一笔业务，再逐步让它面对重试、接管与恢复。实验室和两个地域留下的问题会随实现往前走：每当程序准备再次执行，它凭什么知道，这一次应当继续。




